\documentclass[11pt]{article}

\usepackage[margin=1in]{geometry}
\usepackage{amsmath,amssymb,amsfonts,amsthm}
\usepackage{mathtools}
\usepackage{graphicx}
\graphicspath{{plots/}{rigorous_dyson_theorem/figures/}}
\usepackage{hyperref}
\usepackage{bm}
\usepackage{thmtools}

% Theorem environments
\newtheorem{theorem}{Theorem}[section]
\newtheorem{lemma}[theorem]{Lemma}
\newtheorem{proposition}[theorem]{Proposition}
\newtheorem{corollary}[theorem]{Corollary}
\theoremstyle{definition}
\newtheorem{definition}[theorem]{Definition}
\theoremstyle{remark}
\newtheorem{remark}[theorem]{Remark}
\newtheorem{example}[theorem]{Example}

% Mathematical operators
\DeclareMathOperator{\Tr}{Tr}
\DeclareMathOperator{\bTr}{bTr}
\DeclareMathOperator{\Var}{Var}
\DeclareMathOperator{\supp}{supp}
\newcommand{\E}{\mathbb{E}}
\newcommand{\R}{\mathbb{R}}
\newcommand{\C}{\mathbb{C}}
\newcommand{\N}{\mathbb{N}}
\newcommand{\mcL}{\mathcal{L}}
\newcommand{\mcR}{\mathcal{R}}
\newcommand{\mcX}{\mathcal{X}}
\newcommand{\mcM}{\mathcal{M}}
\newcommand{\mcG}{\mathcal{G}}
\newcommand{\mcS}{\mathcal{S}}
\newcommand{\mcH}{\mathcal{H}}
\newcommand{\mcB}{\mathcal{B}}
\DeclareMathOperator{\tr}{tr}
\newcommand{\norm}[1]{\left\|#1\right\|}
\newcommand{\abs}[1]{\left|#1\right|}
\newcommand{\bigboxplus}{\mathop{\boxplus}}

\title{Orthogonal Hyper-Connections:\\Why the Birkhoff Polytope Breaks Dynamical Isometry}
\author{Leo Brunswic}
\date{}

\begin{document}
\maketitle
\begin{abstract}
Hyper-Connections replace the identity skip connection in residual networks with a learnable mixing matrix $A_q \otimes I_p$, yielding richer information routing and substantial performance gains. The dominant stabilisation strategy constrains $A_q$ to the Birkhoff polytope of doubly stochastic matrices, motivated by mass preservation. We show this intuition is flawed. Using an operator-valued Dyson equation that collapses the spectral problem from network width $N$ to the small twist dimension $q$, we compute the exact limiting Jacobian spectrum for deep networks with structured skip connections. Doubly stochastic matrices cause \emph{partial spectral collapse}: eigenvalues inside the unit disk compound across layers, and eigenspace misalignment accelerates the decay, reintroducing the vanishing-gradient pathology. Orthogonal skip connections eliminate both failure modes---unit-circle eigenvalues prevent contraction, and group closure prevents misalignment---while providing a strictly richer function class. We validate the spectral predictions against Monte Carlo simulation and confirm the practical consequence on a modified Tiny Recursive Model evaluated on the ARC-AGI-1 benchmark: orthogonal skip connections converge faster and reach higher accuracy than their bistochastic counterparts.
\end{abstract}

% \tableofcontents
%!TEX root = ../../main.tex

\section{Introduction}\label{sec:intro}

The residual connection~\cite{he2016deep}---the per-layer update $x^{l+1} = F(x^l) + x^l$---is the dominant connectivity pattern in modern deep learning, underpinning Transformers~\cite{vaswani2017attention} and virtually every large-scale architecture deployed today. Its variants---Pre-Norm, DeepNorm~\cite{wang2022deepnorm}---pushed to thousands of layers. Yet every variant keeps the skip connection as a scalar multiple of the identity, anchoring the neural function to identity and bounding its expressivity.

Hyper-Connections (HC)~\cite{zhu2024hyperconnections} broke out of this constraint by replacing the identity with a learnable mixing matrix $A_q \in \R^{q \times q}$ acting on $q$ parallel residual streams:
\begin{equation}\label{eq:hc-skip}
x^{l+1} = \phi(W^l x^l) + (A_q^l \otimes I_p)\, x^l, \quad N = qp.
\end{equation}
The network learns \emph{which} information flows where---a strictly richer connection pattern. The gains are most dramatic in the Mixture-of-Experts (MoE) setting, where HC halved the training tokens needed to match baseline on OLMoE and improved BBH and GSM8K by $+7$ points on DeepSeek's 27B MoE model~\cite{zhu2024hyperconnections,xie2025mhc}. Both MoE and HC are learnable routing mechanisms---one for tokens across experts, the other for residual streams across layers---and both face the same stability challenge: unconstrained, they diverge (signal gains exceeding $3000\times$ at 27B scale~\cite{xie2025mhc}).

Manifold-Constrained Hyper-Connections (mHC)~\cite{xie2025mhc} tames this instability by projecting $A_q$ onto the Birkhoff polytope of doubly stochastic matrices via the Sinkhorn-Knopp algorithm. The intuition is natural: doubly stochastic matrices preserve total mass and act as transport plans \cite{villani2021topics}, so they should maintain proper information flow. Concurrent work explores alternative parametrisations of the same polytope~\cite{kromhc2026,noguer2026}, but none performs a spectral analysis of the resulting Jacobian. The doubly stochastic constraint is justified by mass preservation and empirical stability---not by any study of spectral structure.

This is where the argument breaks down.\label{sec:spectral-failure} Training deep networks requires that the singular values of the input-output Jacobian $J = \prod_{l=1}^{L} Y^l$ remain concentrated near~$1$---a property called \emph{dynamical isometry}~\cite{saxe2014exact,pennington2017resurrecting}. Without it, gradients either explode or vanish exponentially. Tarnowski et al.~\cite{tarnowski2019dynamical} proved that for scalar skip connections, dynamical isometry is universal: for any activation function, it is achieved when a single condition on the weight variance is met. A generalization of their free probability method to $(A_q^l \otimes I_p)$-twisting of skip-connection seems within reach.  

We extend the theory to the operator-valued setting via operator-valued free probability~\cite{voiculescu1995operations,dykema2006multilinear}, where the Kronecker structure of~\eqref{eq:hc-skip} collapses the spectral problem from network width~$N$ to twist dimension~$q$. This reveals two failure modes of doubly stochastic skip connections. The first is \emph{eigenvalue contraction}: a doubly stochastic matrix has its Perron eigenvalue pinned at one, but generically all others lie strictly inside the unit disk, and deep composition drives $|\lambda|^L \to 0$. The second is \emph{eigenspace misalignment}: eigenbases of successive layers are unrelated, so composition scrambles directions and accelerates the collapse beyond what per-layer spectra predict. Together, these produce a partial spectral collapse---a growing fraction of the Jacobian's singular values drifting toward zero---that no reparametrisation of the Birkhoff polytope can escape (Appendix~\ref{app:spectral-stalling}).

Orthogonal matrices eliminate both failure modes: all eigenvalues lie on the unit circle, so no contraction is possible, and group closure under composition prevents misalignment at any depth. We propose replacing the Birkhoff constraint with the orthogonal group $O(q)$, parametrised via the Cayley transform~\cite{li2020cayley}. Beyond spectral preservation, this provides a strictly richer function class (the linear span of $O(q)$ is the full algebra $M_q(\R)$, dimension $q^2$, versus $(q{-}1)^2 + 1$ for the Birkhoff polytope) and implicit geometric nonlinearity from the curvature of the orthogonal manifold.

\medskip
\noindent\textbf{Contributions.}
\begin{enumerate}
\item \textbf{Spectral diagnosis.} We identify eigenvalue contraction and eigenspace misalignment as the mechanisms by which doubly stochastic skip connections break dynamical isometry, and show that this collapse converts to concrete capacity loss in modern training (\emph{spectral stalling}).

\item \textbf{Orthogonal prescription.} We propose orthogonal skip connections as a principled alternative that preserves dynamical isometry at all depths, with a Cayley-transform parametrisation for efficient Riemannian optimisation~\cite{li2020cayley,lezcano2019trivializations}.

\item \textbf{Implicit differentiation} We optimize both Sinkhorn and Cayley fixed-point methods via implicit differentiation based upon \cite{A unified framework for implicit sinkhorn differentiation}

\item \textbf{Operator-valued Dyson pipeline.} We develop the first numerical implementation of the full operator-valued free probability pipeline---from the matrix Dyson equation through Dykema's twisted S-transform multiplicativity to multi-layer spectral densities.

\item \textbf{Experimental validation.} We confirm the spectral predictions against Monte Carlo simulation and validate the practical consequence on a modified Tiny Recursive Model~\cite{jolicoeur2025trm} evaluated on ARC-AGI-1~\cite{chollet2019arc}: orthogonal skip connections converge faster and reach higher accuracy than bistochastic ones.
\end{enumerate}

%!TEX root = ../../main.tex

\section{Related Work}\label{sec:related-work}

\subsection{Hyper-Connections and structured skip connections}

The standard residual connection $x^{l+1} = F(x^l) + x^l$~\cite{he2016deep} has been extended in several directions. Hyper-Connections (HC)~\cite{zhu2024hyperconnections} replace the identity skip with a learnable mixing matrix $A_q \otimes I_p$ acting on $q$ residual streams, yielding substantial performance gains on language modeling. Manifold-Constrained Hyper-Connections (mHC)~\cite{xie2025mhc} address the resulting training instability by projecting $A_q$ onto the Birkhoff polytope via the Sinkhorn-Knopp algorithm. Two concurrent works explore related directions: KromHC~\cite{kromhc2026} parametrises doubly stochastic matrices as Kronecker products of smaller bistochastic factors via the Birkhoff-von Neumann theorem, and Noguer I Alonso~\cite{noguer2026} proposes an operator-constrained framework. All of these remain within the Birkhoff polytope. Our work departs from this line by showing that the polytope's contractive geometry causes spectral collapse, and proposes the orthogonal group as the correct constraint manifold.

\subsection{Orthogonal constraints in neural networks}

Orthogonal and unitary weight constraints have a rich history in recurrent networks, where they prevent gradient decay across time steps~\cite{arjovsky2016unitary,wisdom2016full}. Lezcano-Casado~\cite{lezcano2019trivializations} developed a general framework for gradient-based optimisation on matrix manifolds via trivializations (exponential map, Cayley transform), enabling efficient training with hard orthogonal constraints. We apply these parametrisation techniques not to weight matrices but to the skip connection matrices $A_q$ in Hyper-Connections, motivated by our spectral analysis showing that orthogonality of $A_q$ is the decisive property for preserving dynamical isometry.

\subsection{Signal propagation and mean-field theory}

The mean-field theory of deep networks~\cite{schoenholz2017deep} studies the propagation of pre-activation moments, identifying the edge-of-chaos phase transition. Yang and Schoenholz~\cite{yang2017mean} extended this to residual networks. Pennington, Schoenholz, and Ganguli~\cite{pennington2017resurrecting} connected signal propagation to the Jacobian's singular value distribution, introducing dynamical isometry for nonlinear networks. These works characterise signal propagation through scalar order parameters; our work computes the full spectral density, revealing failure modes---such as partial spectral collapse in bistochastic skip connections---that are invisible to mean-field analysis.

\subsection{Random matrix theory for neural networks}

Pennington and Worah~\cite{pennington2017nonlinear} developed nonlinear random matrix theory for single-layer networks using the S-transform. Xiao et al.~\cite{xiao2018dynamical} extended this to convolutional architectures. Tarnowski et al.~\cite{tarnowski2019dynamical} derived the universal spectral equation for deep residual networks, showing the Jacobian spectrum depends on the activation only through the effective cumulant $c$, with dynamical isometry at $c = 1/2$. We build on~\cite{tarnowski2019dynamical} in two ways: providing the numerical pipeline needed to solve the spectral equation, and extending the theory from scalar skip connections to the Kronecker structure arising in Hyper-Connections.

\subsection{Operator-valued free probability}

Operator-valued free probability was introduced by Voiculescu~\cite{voiculescu1995operations} and developed combinatorially by Speicher~\cite{speicher1998combinatorial}. The analytic theory was advanced by Belinschi, Mai, and Speicher~\cite{belinschi2017analytic}. The quadratic matrix equation in our Dyson solver was studied by Helton, Far, and Speicher~\cite{helton2007operator}. The twisted multiplicativity of the operator-valued S-transform (Theorem~\ref{thm:dykema}) is due to Dykema~\cite{dykema2006multilinear}; see~\cite{nica2006lectures,mingo2017free} for comprehensive treatments. Our contribution is computational: the first numerical implementation of the full operator-valued pipeline, specialised to the Kronecker structure arising in Hyper-Connections.

%!TEX root = ../../main.tex

\section{Spectral Theory and Computation}\label{sec:spectral-theory}

We now develop the spectral machinery that makes the bistochastic-vs-orthogonal comparison quantitatively precise: the scalar Dyson equation for standard residual networks~\cite{tarnowski2019dynamical}, its operator-valued extension for Kronecker-structured skip connections, and the numerical pipeline that solves both. Full derivations are deferred to the appendices.

\subsection{From scalar Dyson equation to dynamical isometry}\label{sec:dyson-resnet}

A residual network with $L$ layers of width $N$ has forward propagation $x^{l+1} = \phi(W^l x^l) + a\, x^l$, producing the input-output Jacobian
\begin{equation}\label{eq:jacobian}
J = \prod_{l=1}^{L} Y^l, \qquad Y^l = D^l W^l + aI,
\end{equation}
where $D^l = \mathrm{diag}(\phi'(h^l_i))$ is the diagonal matrix of activation derivatives. In the mean-field limit ($N \to \infty$), $D^l$ concentrates and becomes effectively deterministic, so that $D^l W^l$ inherits the isotropy of $W^l$~\cite{pennington2017resurrecting}.

For standard residual networks with scalar skip $a$, a Schwinger--Dyson calculation (Appendix~\ref{app:scalar-recovery}) yields a scalar self-consistent equation---the \emph{scalar Dyson equation}---for the order parameter $m(z)$. For $L$ identical layers, free multiplicative convolution gives a closed-form \emph{$z_1$-mapping} that converts the single-layer solution to the $L$-layer Cauchy transform without numerically fragile S-transform inversion. Together, these recover the \emph{universal spectral equation} of Tarnowski et al.~\cite{tarnowski2019dynamical}: the Jacobian spectrum depends on the activation function $\phi$ only through the effective cumulant $c = L\sigma^2$, where $\sigma^2 = \sigma_w^2\,\E[\phi'(\sqrt{q}\,Z)^2]$. A network achieves \emph{dynamical isometry}\label{def:dynamical-isometry}---singular values of its Jacobian concentrated near $1$---when $c = 1/2$, regardless of activation function or depth.

This is the scalar theory; it applies when $A_q = (a)$---a single scalar skip. Hyper-Connections use a $q \times q$ twist matrix, and here the scalar theory breaks down: the scalar trace $\frac{1}{N}\Tr$ averages over $A_q$'s distinct spectral sectors, erasing the information needed to distinguish bistochastic from orthogonal skip connections. The operator-valued Dyson equation (Proposition~\ref{prop:kronecker-collapse}) resolves individual sectors by promoting the order parameter from a scalar $m(z) \in \C$ to a matrix $M(z) \in M_q(\C)$.

\subsection{Kronecker collapse to $q \times q$}\label{sec:kronecker-skip}

Hyper-Connections~\cite{zhu2024hyperconnections} replace the scalar skip with a Kronecker product $A = A_q \otimes I_p$ ($N = qp$), where $A_q \in \R^{q \times q}$ is the twist matrix encoding inter-stream mixing. The Jacobian factor becomes $Y^l = (A_q^l \otimes I_p) + D^l W^l$.

When $A_q$ has non-trivial eigenvalue structure, the scalar trace $(1/N)\Tr$ averages over $A_q$'s spectral sectors and loses the block geometry. The correct tool is the \emph{operator-valued Dyson equation} over the base algebra $\mcB = M_q(\C)$~\cite{voiculescu1995operations,speicher1998combinatorial}. The Kronecker structure respects the block decomposition of $M_N(\C)$ into $q \times q$ blocks of size $p \times p$, and the isotropy of $W$ ensures the self-energy has no off-diagonal $p \times p$ structure (Appendix~\ref{app:kronecker-derivation}). The result is:

\begin{proposition}[Kronecker collapse]\label{prop:kronecker-collapse}
Under the Kronecker structure and isotropy, the $\mcB$-valued order parameter $M(z) \in M_q(\C)$ satisfies
\begin{equation}\label{eq:matrix-dyson}
M(z) = A_h(M) \cdot \bigl(zI_q - A_h(M)^\top A_h(M)\bigr)^{-1},
\end{equation}
where $A_h(M) := A_q + \sigma^2 M$ is the $q \times q$ dressed matrix. The scalar Cauchy transform is $G(z) = \frac{1}{q}\Tr_q\bigl(zI_q - A_h(M)^\top A_h(M)\bigr)^{-1}$. When $q = 1$, the scalar Dyson equation~\eqref{eq:dyson-fp} is recovered.
\end{proposition}

\noindent This collapses the spectral problem from $O(N^3)$ to $O(q^6)$, making spectral analysis of Hyper-Connections tractable at any network width.

\subsection{Numerical pipeline}\label{sec:numerical}

\noindent\textbf{Scalar solver.}
We solve~\eqref{eq:dyson-fp} by Newton's method with analytical derivatives, sweeping the $z$-grid from large $|z|$ to small $|z|$ so each solve is seeded from the previous solution (\emph{branch continuation}). Full details are in Appendix~\ref{app:numerical-details}.

\medskip
\noindent\textbf{Multi-layer: the $z_1$-mapping.}\label{sec:z1-mapping}
For $L$ identical layers, the S-transform multiplicative property $S_{J^\top J}(w) = S_{Y^\top Y}(w)^L$ reduces the $L$-layer problem to a reparametrisation of the single-layer Cauchy transform, avoiding numerically fragile S-transform inversion (Appendix~\ref{app:scalar-recovery}).

\medskip
\noindent\textbf{Matrix and operator-valued solvers.}\label{sec:matrix-dyson-numerics}
For Kronecker skip connections, the $q \times q$ matrix equation~\eqref{eq:matrix-dyson} is solved by Newton's method in $\C^{q^2}$ at cost $O(q^6)$ per step. For heterogeneous multi-layer networks ($A_q^{l_1} \neq A_q^{l_2}$), the full operator-valued pipeline uses Dykema's twisted S-transform multiplicativity~\cite{dykema2006multilinear} (Appendix~\ref{app:ov-s-transform}) at cost $O(Lq^{10})$ per $z$-point; in practice $q \leq 4$ and $O(1)$ outer iterations suffice with continuation (Appendix~\ref{app:ov-pipeline}).

% Backward-compatibility labels for old section structure
\label{sec:num-dyson}\label{sec:ov-free-prob}\label{sec:matrix-dyson-theory}%
\label{sec:ov-multilayer}\label{sec:identical-vs-hetero}%
\label{sec:w-domain-newton}\label{sec:ov-pipeline-numerics}

%!TEX root = ../../main.tex

\section{Experiments}\label{sec:experiments}

We validate the theoretical predictions in two ways: spectral density computations confirm that bistochastic skip connections produce the predicted partial collapse while orthogonal ones preserve dynamical isometry; training experiments on a modified Tiny Recursive Model confirm the practical consequence.

\subsection{Bistochastic vs.\ orthogonal skip connections}\label{sec:bistochastic-vs-ortho}

The Kronecker framework of Section~\ref{sec:kronecker-skip} enables a direct spectral comparison between the two constraint manifolds for $A_q$---the central question of this paper.

\medskip
\noindent\textbf{Doubly stochastic $A_q$.}
A doubly stochastic matrix has Perron eigenvalue $\lambda_1 = 1$ and remaining eigenvalues $|\lambda_j| < 1$ generically. The operator-valued Dyson equation predicts that the Perron sector contributes a peak near $\sigma = 1$, while contracting sectors produce mass near $\sigma = 0$. Across $L$ layers, contraction compounds: $|\lambda_j|^{2L} \to 0$. This is the \emph{partial spectral collapse} identified in Section~\ref{sec:spectral-failure}.

\medskip
\noindent\textbf{Orthogonal $A_q$.}
An orthogonal matrix has $|\lambda_j| = 1$ for all $j$. No sector contracts. The group structure ensures composition is orthogonal: no misalignment, no drift, at any depth.

\subsection{Spectral density comparison}\label{sec:spectral-experiments}

\noindent\textbf{Prediction.}
The scalar Dyson equation depends on $A_q$ only through $\|A_q\|_F^2/q$, so it cannot distinguish between twist matrices with identical Frobenius norms but different eigenvalue structures. For identity and orthogonal $A_q$ (both with $\|A_q\|_F^2/q = 1$), the scalar prediction should match Monte Carlo. For bistochastic $A_q$ (also $\|A_q\|_F^2/q = 1$ but with contracting eigenvalues), we predict a discrepancy: the scalar theory misses the partial spectral collapse caused by the contracting sectors. For Gaussian $A_q$ ($\|A_q\|_F^2/q < 1$), the scalar theory predicts a shifted density, but Monte Carlo should show complete collapse.

\medskip
\noindent\textbf{Measurement.}
We compare the scalar Dyson prediction against Monte Carlo simulation for four classes of twist matrices at $q = 4$, $p = 200$: identity, random doubly stochastic, Haar-random orthogonal, and Gaussian $A_q = G/\sqrt{q}$. Each class is evaluated at depths $L \in \{10, 20, 30\}$ with ReLU activation and $c_2 L = 0.05$ (500 Monte Carlo samples per configuration).

\medskip
\noindent\textbf{Interpretation.}
Figure~\ref{fig:scalar-vs-mc-panel} confirms all four predictions. For identity and orthogonal $A_q$, the Monte Carlo histograms match the scalar Dyson prediction across all depths---unit-circle eigenvalues preserve dynamical isometry and the scalar theory suffices. For bistochastic $A_q$, the three contracting eigenvalue sectors ($|\lambda_j| < 1$) send $75\%$ of the singular value mass into a Dirac spike at $\sigma = 0$ that is invisible on the plotted $y$-range; the scalar theory, blind to eigenvalue structure, misses this collapse entirely. The Gaussian case ($\|A_q\|_F^2/q = 0.60 < 1$) is even more catastrophic: all Monte Carlo mass collapses to zero. The systematic discrepancy between scalar prediction and Monte Carlo for non-isotropic $A_q$ is precisely what motivates the operator-valued Dyson equation of Section~\ref{sec:kronecker-skip}.



% ── Scalar theory vs Monte Carlo panel (4 matrix types × 3 depths) ──
% Generated by experiments/panel_sv_density_4matrices.py on GPU (RTX 4090).
% Scalar theory uses identity prediction for orthogonal and bistochastic
% columns, since the scalar Dyson equation depends only on ||A||_F^2/q.
\begin{figure}[t]
\centering
\includegraphics[width=\textwidth]{panel_sv_4matrices_q4_p200_n500.pdf}
\caption{Scalar theory vs.\ Monte Carlo singular value density for four skip-connection types ($q=4$, $p=200$, $c_2 L = 0.05$, $\eta = 0.02$, $n_{\mathrm{samples}} = 500$). Rows correspond to depths $L \in \{10, 20, 30\}$; columns to $A_q \in \{I_q,\, \text{bistochastic},\, \text{orthogonal},\, G/\sqrt{q}\}$. The scalar theory (blue) is computed from the identity reduction for the first three columns (since $\|A_q\|_F^2/q = 1$ for identity and orthogonal, and the scalar equation ignores eigenvalue structure). \textbf{Key observations:} (i)~Identity and orthogonal Monte Carlo match the scalar prediction; (ii)~bistochastic Monte Carlo concentrates $75\%$ of its mass in a Dirac spike at $\sigma = 0$ (three contracting sectors with $\sigma_k^L \to 0$), invisible on the plotted $y$-range; (iii)~Gaussian $A_q$ ($\|A_q\|_F^2/q = 0.60 < 1$) sends all MC mass to zero. The discrepancy between scalar theory and MC for non-isotropic $A_q$ motivates the operator-valued treatment.}
\label{fig:scalar-vs-mc-panel}
\end{figure}

\subsection{Training experiments: TRM on ARC-AGI-1}\label{sec:trm-experiments}

\noindent\textbf{Prediction.}
The spectral collapse in bistochastic skip connections should manifest as slower convergence and reduced accuracy, particularly in architectures that apply their layers recursively (amplifying effective depth). Orthogonal skip connections, which preserve dynamical isometry, should converge faster and achieve higher accuracy.

\medskip
\noindent\textbf{Setup.}
We modify the Tiny Recursive Model (TRM)~\cite{jolicoeur2025trm}, a minimalist recursive architecture for abstract reasoning with $\sim\!7$M parameters. TRM applies its two-layer network recursively, making it exquisitely sensitive to skip-connection geometry: each recursive application compounds the Jacobian product, amplifying any spectral pathology. We replace TRM's standard skip connections with Kronecker-structured connections $A_q \otimes I_p$ and compare three constraint manifolds for $A_q$:
\begin{itemize}
\item \textbf{Doubly stochastic} (Sinkhorn-Knopp projection, as in mHC~\cite{xie2025mhc});
\item \textbf{Orthogonal} (exponential map parametrisation~\cite{lezcano2019trivializations});
\item \textbf{Unconstrained} (baseline, no manifold constraint).
\end{itemize}
All models are trained on the ARC-AGI-1 benchmark~\cite{chollet2019arc} with identical hyperparameters except for the skip-connection constraint.

\medskip
\noindent\textbf{Results.}
% TODO: Insert table/figure with TRM results once experiments are complete.
% Expected: orthogonal > bistochastic > unconstrained in accuracy;
% orthogonal converges in fewer epochs; unconstrained diverges at higher recursion depths.

\medskip
\noindent\textbf{Interpretation.}
The results are consistent with the spectral theory: orthogonal skip connections preserve gradient flow across recursive applications, avoiding the capacity loss from partial spectral collapse. The advantage over bistochastic is most pronounced at higher recursion depths, exactly where the compounding eigenvalue contraction becomes most severe.

%!TEX root = ../../main.tex

\section{Discussion}\label{sec:discussion}

The Birkhoff polytope is the wrong constraint manifold for Hyper-Connections. Its contractive eigenvalue geometry causes partial spectral collapse that worsens with depth, and eigenspace misalignment between layers accelerates the decay. Orthogonal skip connections eliminate both failure modes while providing a strictly richer function class (dimension $q^2$ vs.\ $(q-1)^2 + 1$). The operator-valued Dyson pipeline developed here---collapsing the $N$-dimensional spectral problem to dimension $q$ via Kronecker structure---is what makes this comparison quantitatively precise.

\medskip
\noindent\textbf{Spectral stalling.}
Partial spectral collapse is not merely an asymptotic concern. Four independent mechanisms in modern training each impose a hard cutoff $\sigma^*$ on the Jacobian singular values, below which gradient directions are operationally invisible to the optimiser (Appendix~\ref{app:spectral-stalling}). The binding constraint is $\sigma^* = \max(\sigma^*_{\mathrm{time}}, \sigma^*_{\mathrm{Adam}}, \sigma^*_{\mathrm{bf16}}, \sigma^*_{\mathrm{noise}})$, arising respectively from finite training horizon~\cite{saxe2014exact}, Adam's $\epsilon$-normalisation~\cite{kingma2015adam}, bf16 swamping~\cite{micikevicius2018mixed}, and mini-batch noise. In typical modern training (Adam, bf16, $T \sim 10^5$ steps), the bf16 floor and finite-time horizon are the dominant cutoffs. Our Dyson pipeline gives the spectral measure $\mu_\Sigma$ as a function of $A_q$, so the \emph{stalling fraction}---the mass below $\sigma^*$---is computable: $\rho_{\mathrm{stall}}(A_q, L, \sigma^*) = \int_0^{\sigma^*} d\mu_\Sigma(\sigma)$. For doubly stochastic $A_q$, contracting eigenvalue sectors push $\rho_{\mathrm{stall}}$ upward with depth. For orthogonal $A_q$, no sector contracts, and $\rho_{\mathrm{stall}}$ remains small. The effective training dimension is $n_{\mathrm{eff}} = n \cdot (1 - \rho_{\mathrm{stall}})$: spectral collapse becomes \emph{capacity loss}.

\medskip
\noindent\textbf{Limitations.}
The Dyson equation requires isotropy of $D^l W^l$ (Section~\ref{sec:dyson-resnet}), which holds for i.i.d.\ Gaussian or orthogonally invariant weights but excludes sparse, low-rank, convolutional, or batch-normalised weight matrices. The mean-field treatment of the activation derivatives $D^l$ becomes exact only as $N \to \infty$~\cite{pennington2017resurrecting}. The Kronecker structure $A_q \otimes I_p$ captures multi-stream architectures but not general structured skip connections. The exponential-map parametrisation of $O(q)$ costs a matrix exponential per forward pass---negligible for the small $q \leq 8$ relevant in practice, but requiring more efficient parametrisations at larger $q$.

\medskip
\noindent\textbf{Outlook.}
The most immediate next step is testing whether the spectral advantage persists at scale, in the large language models where mHC was originally developed. Extending the spectral analysis from initialisation to training dynamics---tracking how the Jacobian spectrum evolves as isotropy progressively breaks down---would connect the present theory to the dynamics of learning.


\appendix
%!TEX root = ../../main.tex

\section{Free Probability Tools}\label{app:free-prob}

This appendix provides the basic framework of free probability theory necessary for understanding the spectral analysis of Jacobian products in deep residual neural networks. We introduce the R-transform and S-transform, which play analogous roles to the moment and cumulant generating functions in classical probability, and explain their application to free multiplicative convolution. For comprehensive treatments, see Nica and Speicher~\cite{nica2006lectures}, Mingo and Speicher~\cite{mingo2017free}, and Hiai and Petz~\cite{hiai2000semicircle}.

\subsection{R-transform and free cumulants}

The R-transform is a fundamental tool in free probability that encodes the free cumulants of a probability distribution.

\begin{definition}[Free Cumulants and R-transform]\label{def:r-transform}
Let $\mu$ be a probability measure on $\R$ with moments $m_n = \int x^n d\mu(x)$ for $n \geq 1$. The \emph{free cumulants} $(\kappa_n)_{n \geq 1}$ of $\mu$ are defined combinatorially via non-crossing partitions (see \cite{nica2006lectures}, Lecture 11).

The \emph{R-transform} of $\mu$ is the formal power series
\begin{equation}\label{eq:r-transform-def}
\mathcal{R}_\mu(z) = \sum_{n=1}^\infty \kappa_n z^n.
\end{equation}
\end{definition}

The R-transform is related to the Cauchy (Stieltjes) transform via a functional equation.

\begin{proposition}[R-transform functional equation]\label{prop:r-cauchy-relation}
Let $G_\mu(z) = \int \frac{d\mu(\lambda)}{z - \lambda}$ be the Cauchy transform of $\mu$. Then the R-transform satisfies
\begin{equation}\label{eq:r-functional}
\mathcal{R}_\mu(G_\mu(z)) + \frac{1}{G_\mu(z)} = z.
\end{equation}

Equivalently, if $K_\mu(z) = 1/G_\mu(1/z)$ denotes the inverse Cauchy transform, then
\begin{equation}
\mathcal{R}_\mu(z) = K_\mu(z) - \frac{1}{z}.
\end{equation}
\end{proposition}

\begin{proof}[Proof sketch]
The relation \eqref{eq:r-functional} follows from the combinatorial definition of free cumulants and the moment-cumulant formula (Theorem 11.12 in \cite{nica2006lectures}). For a complete proof, see \cite{nica2006lectures}, Lecture 12.
\end{proof}

\begin{example}[Free Poisson distribution]\label{ex:free-poisson}
The free Poisson distribution with rate $\lambda > 0$ and jump size $\alpha \in \R$ has R-transform
\begin{equation}
\mathcal{R}(z) = \frac{\lambda \alpha}{1 - \alpha z}.
\end{equation}

This corresponds to the limiting spectral distribution of the matrix $\lambda^{1/2} X$, where $X$ has i.i.d.\ entries with variance $1/N$; it is the free analogue of the classical Marchenko--Pastur law~\cite{marchenko1967distribution} (see \cite{nica2006lectures}, Proposition 12.11).
\end{example}

\subsection{S-transform and multiplicative convolution}

While the R-transform linearizes free \emph{additive} convolution, the S-transform linearizes free \emph{multiplicative} convolution, making it the key tool for analyzing products of free random matrices.

\begin{definition}[S-transform]\label{def:s-transform}
Let $\mu$ be a probability measure on $\R^+$ with Cauchy transform $G_\mu(z) = \int \frac{d\mu(\lambda)}{z - \lambda}$. Define the \emph{moment function} $\Psi_\mu(z) := z\,G_\mu(z) - 1$, and let $\chi_\mu$ denote its functional inverse: $\Psi_\mu(\chi_\mu(w)) = w$. The \emph{S-transform} of $\mu$ is
\begin{equation}\label{eq:s-transform-def}
S_\mu(w) = \frac{w+1}{w \cdot \chi_\mu(w)}.
\end{equation}
Equivalently, for each $w$ in a neighborhood of~$0$, let $z_w = \chi_\mu(w)$ be the unique solution of $z\,G_\mu(z) - 1 = w$. Then
\begin{equation}\label{eq:s-transform-via-G}
S_\mu(w) = \frac{G_\mu(z_w)}{w}.
\end{equation}
\end{definition}

\begin{remark}[Reconciliation with the Cauchy--S-transform relation]\label{rem:reconciliation_cauchy__s_transform_relation}
The familiar identity
\[
\frac{G(z)}{zG(z)-1} = S\bigl(zG(z)-1\bigr),
\]
is obtained by setting $w = zG(z) - 1 = \Psi(z)$ in~\eqref{eq:s-transform-via-G}. Indeed, $z_w = z$ solves $z\,G(z) - 1 = w$ by construction, so $S(w) = G(z)/w = G(z)/(zG(z)-1)$.
\end{remark}

\begin{remark}\label{rem:s_transform_particularly_natural_measures}
The S-transform is particularly natural for measures on $\R^+$, which arise as the spectral distributions of $XX^T$ for random matrices $X$. For the extension to general signed measures (relevant for eigenvalues as opposed to singular values), see \cite{nica2006lectures}, Lecture 18.
\end{remark}

The key property of the S-transform is its linearization of free multiplicative convolution.

\begin{theorem}[S-transform linearization]\label{thm:s-linearization}
Let $a, b$ be freely independent random variables in a $C^*$-probability space with spectral measures $\mu_a, \mu_b$ supported on $\R^+ \cup \{0\}$. Then the spectral measure $\mu_{ab}$ of the product $ab$ satisfies
\begin{equation}\label{eq:s-multiplicative}
S_{\mu_{ab}}(z) = S_{\mu_a}(z) \cdot S_{\mu_b}(z).
\end{equation}

More generally, for a product of $L$ freely independent random variables $a_1, \ldots, a_L$,
\begin{equation}
S_{\mu_{a_1 \cdots a_L}}(z) = \prod_{l=1}^L S_{\mu_{a_l}}(z).
\end{equation}
\end{theorem}

\begin{proof}
See \cite{nica2006lectures}, Lecture 18, or \cite{voiculescu1992free} for the complete proof. The key idea is that the S-transform encodes the ``compression'' of the spectral measure under multiplication, and free independence ensures that compressions compose multiplicatively.
\end{proof}

\subsection{Free multiplicative convolution and asymptotic freeness}

\begin{definition}[Free multiplicative convolution]\label{def:free-multiplicative-convolution}
Let $\mu, \nu$ be probability measures on $\R^+$. The \emph{free multiplicative convolution} $\mu \boxtimes \nu$ is the unique probability measure on $\R^+$ whose S-transform satisfies
\begin{equation}
S_{\mu \boxtimes \nu}(z) = S_\mu(z) \cdot S_\nu(z).
\end{equation}
\end{definition}

The relevance of free multiplicative convolution to random matrix theory comes from the following fundamental result.

\begin{theorem}[Asymptotic freeness of independent matrices]\label{thm:asymptotic-freeness}
Let $A_N, B_N \in \R^{N \times N}$ be sequences of random matrices such that:
\begin{enumerate}
\item $A_N, B_N$ are independent for each $N$;
\item The empirical spectral distributions $\mu_{A_N}$ converge weakly to deterministic limits $\mu_A, \mu_B$ as $N \to \infty$;
\item Standard concentration and moment conditions hold.
\end{enumerate}

Then the empirical spectral distribution of $A_N B_N$ converges weakly to $\mu_A \boxtimes \mu_B$ as $N \to \infty$.
\end{theorem}

\begin{proof}[Proof sketch]
The key observation is that independence of matrices implies asymptotic freeness of their ``traces'' in the large-$N$ limit. This is a consequence of concentration of measure: fluctuations in $\frac{1}{N}\Tr(A_N^k B_N^l \cdots)$ become negligible, leaving only the ``free'' contributions. For a rigorous proof using subordination functions, see \cite{anderson2010introduction}, Theorem 5.4.5.
\end{proof}

\begin{corollary}[Jacobian products in ResNets]\label{cor:jacobian-free-convolution}
For a residual network with Jacobian $J = \prod_{l=1}^L Y^l$ where each $Y^l = aI + D^l W^l$ and the $W^l$ are independent random matrices, the limiting spectral measure of $J^T J$ satisfies
\begin{equation}\label{eq:jacobian-spectral}
\mu_{J^T J} = \bigboxplus_{l=1}^L \mu_{(Y^l)^T Y^l},
\end{equation}
where $\bigboxplus$ denotes the iterated free multiplicative convolution. Consequently,
\begin{equation}\label{eq:jacobian-s-transform}
S_{J^T J}(z) = \prod_{l=1}^L S_{(Y^l)^T Y^l}(z).
\end{equation}
\end{corollary}

\begin{proof}
Apply Theorem \ref{thm:asymptotic-freeness} inductively. The independence of the weight matrices $W^l$ across layers, combined with concentration of the activation patterns $D^l$, ensures that the factors $Y^l$ become asymptotically free in the large-width limit $N \to \infty$. See Section 5.1 for details.
\end{proof}

\subsection{Connection to Dyson equation and universal spectral equation}

We now connect the S-transform formalism to the Dyson equation framework of Section~\ref{sec:dyson-resnet}, and use this to derive the universal spectral equation for residual networks.

\begin{proposition}[S-transform from Dyson equation]\label{prop:s-from-dyson}
Let $Y = A + X$ where $X$ is isotropic with variance $\sigma^2/N$, and assume the deterministic equivalent $M(z)$ is given by the Dyson equation~\eqref{eq:dyson-fp}. Define the order parameters
\begin{equation}
m_{12}(z) = \frac{1}{N}\Tr(M_{12}(z)), \quad m_{21}(z) = \frac{1}{N}\Tr(M_{21}(z)).
\end{equation}

Then the S-transform of the spectral measure of $Y^\top Y$ is obtained as follows. For each $w$ in a neighborhood of $0$, let $z_w$ be the unique solution on the physical (Stieltjes) branch of
\begin{equation}\label{eq:psi-inversion-scalar}
z \, G_1(z) - 1 = w,
\end{equation}
where $G_1(z) = g_{11}(z)$ is the Cauchy transform of $Y^\top Y$ from the Dyson equation. Then
\begin{equation}\label{eq:s-from-order-params}
S_{Y^\top Y}(w) = \frac{G_1(z_w)}{w}.
\end{equation}

For the specific case $Y = aI + X$ with isotropic $X$, explicit formulas are given in \cite{tarnowski2019dynamical}, Section II.
\end{proposition}

\begin{proof}[Proof sketch]
The Cauchy transform $G(z)$ of the singular values of $Y$ is related to the (1,1) block of the resolvent via $G(z) \approx g_{11}(z) = m_{11}(z)$. The S-transform is then obtained by the change of variables in Definition \ref{def:s-transform}. The order parameters $(m_{12}, m_{21})$ encode the full spectral information through the self-consistent Dyson equation. See \cite{burda2010free} for detailed calculations in the case of products of random matrices.
\end{proof}

We can now state the main result connecting free probability to the universal spectral equation.

\begin{theorem}[Universal spectral equation via S-transform]\label{thm:universal-via-s}
Consider a deep residual network with $L$ layers, skip connection parameter $a$, and effective cumulant $c = L\sigma^2$ (Section~\ref{sec:dyson-resnet}). In the large-depth limit $L \to \infty$ with $c^l \sim c/L$ per layer, the S-transform of the Jacobian spectrum satisfies
\begin{equation}\label{eq:s-composition}
S_J(z) = \prod_{l=1}^L S_{Y^l}(z) = a^{2L} \exp\left(\int_0^z \frac{2c}{w(w+1)} dw\right) = a^{2L} \left(\frac{z+1}{z}\right)^{2c}.
\end{equation}

Inverting the S-transform to obtain the Cauchy transform $G(z)$ yields the universal equation
\begin{equation}\label{eq:universal-from-s}
\boxed{a^{2L} G(z) = (zG(z) - 1) \exp\bigl(z(1 - 2cG(z))\bigr),}
\end{equation}
which is exactly the universal spectral equation of Tarnowski et al.~\cite{tarnowski2019dynamical}.
\end{theorem}

\begin{proof}
\textbf{Step 1: S-transform composition.} By Corollary \ref{cor:jacobian-free-convolution},
\[
S_J(z) = \prod_{l=1}^L S_{Y^l}(z).
\]

For $Y^l = aI + D^l W^l$ with small random perturbations, each layer contributes approximately
\[
S_{Y^l}(z) \approx a^2 \exp\left(\frac{2c^l}{z(z+1)}\right),
\]
where $c^l$ is the effective cumulant at layer $l$ (see Proposition \ref{prop:s-from-dyson} and \cite{tarnowski2019dynamical}, Eq. (23)).

\textbf{Step 2: Large-depth limit.} Taking $L \to \infty$ with $c^l = c/L$ yields
\begin{align}
S_J(z) &= \prod_{l=1}^L a^2 \exp\left(\frac{2c/L}{z(z+1)}\right) \\
&= a^{2L} \exp\left(\sum_{l=1}^L \frac{2c/L}{z(z+1)}\right) \\
&= a^{2L} \exp\left(\frac{2c}{z(z+1)}\right) \\
&\to a^{2L} \exp\left(\int_0^z \frac{2c}{w(w+1)} dw\right) \quad \text{(continuum limit)}.
\end{align}

Evaluating the integral:
\begin{align}
\int \frac{2c}{w(w+1)} dw &= 2c \int \left(\frac{1}{w} - \frac{1}{w+1}\right) dw \\
&= 2c \log\left(\frac{w+1}{w}\right) + \text{const}.
\end{align}

Thus,
\[
S_J(z) = a^{2L} \left(\frac{z+1}{z}\right)^{2c}.
\]

\textbf{Step 3: Inversion to Cauchy transform.} We invert the S-transform using the relation $S(w) = G(z_w)/w$ where $z_w$ solves $z\,G(z) - 1 = w$ (Definition~\ref{def:s-transform}). Setting $w = z\,G_L(z) - 1$ gives $S_J(w) = G_L(z)/w$, so
\[
\frac{G_L(z)}{z\,G_L(z) - 1} = a^{2L} \left(\frac{z\,G_L(z)}{z\,G_L(z) - 1}\right)^{2c}.
\]

Let $u := z\,G_L(z)$, so $G_L(z) = u/z$ and the equation becomes
\[
\frac{u/z}{u - 1} = a^{2L} \left(\frac{u}{u - 1}\right)^{2c}.
\]

For the detailed inversion from this implicit equation to the universal form~\eqref{eq:universal-from-s}, one uses the explicit per-layer Dyson structure. Tarnowski et al.~\cite{tarnowski2019dynamical} (Section~II, Eqs.~(18)--(25)) show that the single-layer S-transform for $Y = aI + X$ with isotropic $X$ of variance $c/L$ satisfies $\log S_{Y^\top Y}(w) = 2\log a + 2(c/L)/(w(w+1)) + O(1/L^2)$. Taking the $L$-layer product and the limit $L \to \infty$ (with fixed $c$) yields
\[
a^{2L} G(z) = (zG(z) - 1) \exp\bigl(z(1 - 2cG(z))\bigr),
\]
which is the desired universal spectral equation. The detailed calculation is given in~\cite{tarnowski2019dynamical}, Eq.~(25).
\end{proof}

\begin{remark}[Universality interpretation]\label{rem:universality_interpretation}
Theorem \ref{thm:universal-via-s} shows that the spectral properties of the Jacobian depend \emph{only} on:
\begin{itemize}
\item The depth $L$,
\item The skip connection strength $a$,
\item The effective cumulant $c = \frac{1}{L}\sum_{l=1}^L c^l$, where $c^l$ encodes the variance and activation statistics at layer $l$.
\end{itemize}

The specific choice of activation function $\phi$ (ReLU, tanh, sigmoid, etc.) affects the spectrum \emph{only through} its contribution to $c$. This is the \textbf{spectral universality} observed in \cite{tarnowski2019dynamical}, Figure 3: different activation functions with the same effective cumulant yield identical Jacobian spectra.
\end{remark}

\begin{remark}[Dynamical isometry]\label{rem:dynamical_isometry}
Setting $a = 1$ and solving for the value of $c$ that yields a spectrum concentrated near $\lambda = 1$ (dynamical isometry), one finds $c = 1/2$ (Definition~\ref{def:dynamical-isometry}). This critical value is universal across all activation functions, depending only on the constraint that the spectrum of $J$ should be localized at 1.
\end{remark}

\subsection{Operator-valued S-transform and Dykema's twisted multiplicativity}\label{app:ov-s-transform}

For multi-layer composition with $q > 1$, the S-transform generalises to a $\mcB$-valued map ($\mcB = M_q(\C)$). The $\mcB$-valued S-transform is $S^{(\mcB)}(W) := G^{(\mcB)}(b) \cdot W^{-1}$, where $b$ solves the \emph{Psi-inversion}
\begin{equation}\label{eq:psi-inversion}
b \cdot G^{(\mcB)}(b) - I_q = W.
\end{equation}
This inherits a \emph{twist} from the non-commutativity of $\mcB$:

\begin{theorem}[Dykema, 2006]\label{thm:dykema}
Let $x$ and $y$ be $\mcB$-freely independent positive elements. Then
\begin{equation}\label{eq:twisted-mult}
S^{(\mcB)}_{xy}(W) = S^{(\mcB)}_y(W) \cdot S^{(\mcB)}_x\!\bigl(S^{(\mcB)}_y(W)^{-1} \, W \, S^{(\mcB)}_y(W)\bigr).
\end{equation}
\end{theorem}

\noindent For $L$ layers, this gives the twisted recursion:
\begin{equation}\label{eq:twisted-fold}
S_{\mathrm{cumul}} \leftarrow S^{(\mcB)}_L(W); \quad \text{for } k = L{-}1, \ldots, 1: \;\; W_{\mathrm{tw}} = S_{\mathrm{cumul}}^{-1} W S_{\mathrm{cumul}}, \;\; S_{\mathrm{cumul}} \leftarrow S_{\mathrm{cumul}} \cdot S^{(\mcB)}_k(W_{\mathrm{tw}}).
\end{equation}
When $q = 1$, the twist is trivial (scalars commute) and~\eqref{eq:twisted-fold} reduces to $S_{\mathrm{cumul}} = \prod_l S_l(w)$.

The $L$-layer $\mcB$-valued Cauchy transform is obtained by solving the consistency equation: find $W \in M_q(\C)$ such that
\begin{equation}\label{eq:outer-consistency}
(W + I_q)\bigl(S_{\mathrm{prod}}(W) \cdot W\bigr)^{-1} = z I_q,
\end{equation}
where $S_{\mathrm{prod}}(W)$ is the twisted product~\eqref{eq:twisted-fold}. The scalar density is then $G_L(z) = \frac{1}{qz}\Tr_q(W + I_q)$.

\subsection{Scalar special case: recovery of the classical ResNet result}\label{app:scalar-recovery}

When $q = 1$ (standard residual networks with scalar skip $a$), the operator-valued Dyson equation~\eqref{eq:matrix-dyson} reduces to a scalar fixed-point problem. Together with the $z_1$-mapping for multi-layer composition, this recovers the universal spectral equation of Tarnowski et al.~\cite{tarnowski2019dynamical}.

\begin{proposition}[Scalar Dyson equation and $z_1$-mapping]\label{prop:scalar-dyson}
Let $Y = aI_N + DW$ with effective variance $\sigma^2 := \sigma_w^2\,\E[\phi'(\sqrt{q}\,Z)^2]$, $Z \sim \mathcal{N}(0,1)$.

\noindent\emph{(i) Single layer.} In the limit $N \to \infty$, the order parameter $m(z) := \frac{1}{N}\Tr\!\bigl(Y(zI_N - Y^\top\! Y)^{-1}\bigr)$ satisfies the self-consistent \emph{scalar Dyson equation}\label{app:dyson-derivation}
\begin{equation}\label{eq:dyson-fp}
m(z) = \frac{1}{N}\Tr\bigl(A_h(z)\,(zI_N - A_h(z)^\top A_h(z))^{-1}\bigr), \qquad A_h(z) := aI + \sigma^2 m(z)\, I,
\end{equation}
where $A_h$ is the dressed skip connection. The Cauchy transform is $G(z) = \frac{1}{N}\Tr(zI_N - A_h^\top A_h)^{-1}$ and the spectral density is $\rho(\lambda) = -\pi^{-1}\lim_{\eta \to 0^+}\!\Im\, G(\lambda + i\eta)$.

\noindent\emph{(ii) $L$ identical layers.}\label{prop:z1-mapping} Let $G_1(z)$ be the single-layer Cauchy transform. Then for any $z_1 \in \C^+$:
\begin{equation}\label{eq:z1-to-zL}
z_L(z_1) := z_1 \cdot \left(\frac{z_1 G_1(z_1) - 1}{G_1(z_1)}\right)^{L-1}, \qquad
G_L(z_L) = \frac{z_1 \, G_1(z_1)}{z_L}.
\end{equation}
\end{proposition}

\noindent Part~(i) is derived below via the Schwinger--Dyson approach; part~(ii) is proved in Appendix~\ref{app:z1-proof}. The Jacobian spectrum depends on $\phi$ only through the effective cumulant $c = L\sigma^2$; dynamical isometry is achieved at $c = 1/2$.

\medskip
\noindent\textbf{Derivation of part~(i).}
The argument follows the standard Schwinger--Dyson approach via Gaussian integration by parts (Appendix~\ref{app:gaussian-ibp}).

\medskip
\noindent\textbf{Step 1: Block linearization.}
Let $Y = A + X \in \R^{N \times N}$ where $A = aI_N$ is the deterministic skip and $X$ has i.i.d.\ $\mathcal{N}(0, \sigma^2/N)$ entries (the activation derivative is absorbed into $\sigma^2$ as in Section~\ref{sec:dyson-resnet}). Rather than working directly with the resolvent of $Y^\top Y$, we introduce the $2N \times 2N$ block linearization
\begin{equation}\label{eq:block-lin}
\mcL(z) = \begin{pmatrix} zI_N & -Y \\ -Y^\top & I_N \end{pmatrix},\qquad
\mcG(z) = \mcL(z)^{-1} = \begin{pmatrix} G_{11} & G_{12} \\ G_{21} & G_{22} \end{pmatrix}.
\end{equation}
By the Schur complement, the blocks of $\mcG$ encode the spectral information of $Y$:
\begin{align}
G_{11} &= (zI - YY^\top)^{-1}, \label{eq:G11}\\
G_{22} &= z\,(zI - Y^\top Y)^{-1}, \label{eq:G22}\\
G_{12} &= Y\,(zI - Y^\top Y)^{-1}. \label{eq:G12}
\end{align}
The Cauchy transform of $Y^\top Y$ is $G(z) = \frac{1}{N}\Tr(zI - Y^\top Y)^{-1} = \frac{1}{zN}\Tr G_{22}$, and the off-diagonal order parameter is $m(z) = \frac{1}{N}\Tr G_{12}$.

\medskip
\noindent\textbf{Step 2: Schwinger--Dyson equation.}
Split $\mcL = \mcL_0 + \mcV$ where
\[
\mcL_0 = \begin{pmatrix} zI & -A \\ -A^\top & I \end{pmatrix}, \qquad
\mcV = \begin{pmatrix} 0 & -X \\ -X^\top & 0 \end{pmatrix}.
\]
The identity $\mcL\,\mcG = I_{2N}$ gives $\mcG_0^{-1}\mcG + \mcV\,\mcG = I$, i.e.,
\begin{equation}\label{eq:sd-identity}
\mcG = \mcG_0 - \mcG_0 \,\mcV\, \mcG,
\end{equation}
where $\mcG_0 = \mcL_0^{-1}$ is the ``bare'' Green's function (the resolvent with $X = 0$).

\medskip
\noindent\textbf{Step 3: Wick contraction via Gaussian integration by parts.}
We take $\frac{1}{N}\Tr$ of each $N \times N$ block in~\eqref{eq:sd-identity} and take expectations. The nontrivial terms are those involving $X$ against $\mcG$. We evaluate them using Stein's lemma (Proposition~\ref{prop:matrix-stein}): for each entry $X_{ij} \sim \mathcal{N}(0, \sigma^2/N)$,
\[
\E\bigl[X_{ij}\,(\mcG)_{kl}\bigr] = \frac{\sigma^2}{N}\,\E\!\left[\frac{\partial (\mcG)_{kl}}{\partial X_{ij}}\right].
\]
The resolvent derivative (Lemma~\ref{lem:resolvent-derivative}) gives $\partial \mcG / \partial X_{ij} = \mcG \,(\partial \mcV / \partial X_{ij})\, \mcG$, so the Wick contraction produces terms quadratic in $\mcG$.

Concretely, consider the $(1,2)$ block of~\eqref{eq:sd-identity}. The term involving $X$ is $-\frac{1}{N}\Tr\bigl(\mcG_0^{(11)} X \cdot \mcG^{(22)} + \mcG_0^{(12)} X^\top \mcG^{(12)}\bigr)$, where superscripts denote $N \times N$ blocks. Taking expectations and applying Stein's lemma entry-by-entry:
\begin{align}
\E\!\left[\frac{1}{N}\sum_{i,j} X_{ij}\,(G_{22})_{ji}\right]
&= \frac{\sigma^2}{N^2}\sum_{i,j}\E\!\left[\frac{\partial (G_{22})_{ji}}{\partial X_{ij}}\right] \notag\\
&= \frac{\sigma^2}{N^2}\sum_{i,j}\E\!\left[(G_{12})_{jj}\,(G_{11})_{ii} + \cdots\right] \notag\\
&\;\xrightarrow{N \to \infty}\; \sigma^2\,m(z)\,g_{11}(z), \label{eq:wick-result}
\end{align}
where $g_{11}(z) = \lim_{N\to\infty}\frac{1}{N}\Tr G_{11}$ and $m(z) = \lim_{N\to\infty}\frac{1}{N}\Tr G_{12}$ are the deterministic limits, which exist by concentration (Theorem~\ref{thm:concentration-green}). The factorization in the last step uses the fact that $\frac{1}{N}\Tr(G_{12})$ concentrates around its expectation, so the product of traces converges to the product of limits.

The remaining Wick contractions produce analogous terms involving $m(z)$, $g_{11}(z)$, and $g_{22}(z) = z\,G(z)$.

\medskip
\noindent\textbf{Step 4: Self-consistent closure.}
Inserting the Wick contractions into the block-traced Schwinger--Dyson equation yields a closed system for the $2\times 2$ matrix of order parameters $\bigl(\begin{smallmatrix} g_{11} & m \\ \bar{m} & g_{22}\end{smallmatrix}\bigr)$. The self-energy $\sigma^2 m(z)$ enters as a scalar shift of the skip connection:
\begin{equation}\label{eq:dressed-skip}
A = aI \;\longrightarrow\; A_h(z) := aI + \sigma^2 m(z)\,I = \bigl(a + \sigma^2 m(z)\bigr)\,I.
\end{equation}
The deterministic limit of $\mcG$ equals the Green's function of the dressed block linearization
\[
\bar{\mcG}(z) = \begin{pmatrix} zI & -A_h \\ -A_h^\top & I \end{pmatrix}^{-1}\!,
\]
whose blocks are computed by the Schur complement exactly as in~\eqref{eq:G11}--\eqref{eq:G12} with $Y$ replaced by $A_h$. In particular,
\[
m(z) = \frac{1}{N}\Tr\!\bigl(A_h(zI - A_h^\top A_h)^{-1}\bigr), \qquad G(z) = \frac{1}{N}\Tr(zI - A_h^\top A_h)^{-1},
\]
which is the scalar Dyson equation~\eqref{eq:dyson-fp}. Since $A_h$ depends on $m(z)$, this is a fixed-point equation for $m$ that must be solved numerically (see Appendix~\ref{app:scalar-jacobian}).

\medskip
\noindent\textbf{Step 5: Spectral density recovery.}
Once $m(z)$ (and hence $G(z)$) has been computed on a grid $z_k = \lambda_k + i\eta$ with $\eta > 0$ small, the spectral density of $Y^\top Y$ is obtained by the Stieltjes inversion formula:
\begin{equation}\label{eq:stieltjes-inversion}
\rho(\lambda) = -\frac{1}{\pi}\lim_{\eta \to 0^+}\Im\, G(\lambda + i\eta).
\end{equation}
In practice, $\eta$ is fixed at a small positive value (the broadening parameter) and the density is read off as $\rho(\lambda) \approx -\frac{1}{\pi}\Im\, G(\lambda + i\eta)$. The singular-value density of $Y$ is then $\rho_\sigma(\sigma) = 2\sigma\,\rho(\sigma^2)$ by the change of variables $\lambda = \sigma^2$.

\subsection{Derivation of the Kronecker collapse}\label{app:kronecker-derivation}

We derive the $q \times q$ matrix Dyson equation~\eqref{eq:matrix-dyson} stated in Proposition~\ref{prop:kronecker-collapse}, extending the scalar derivation of Section~\ref{app:scalar-recovery} to the operator-valued setting.

\medskip
\noindent\textbf{Step 1: Block decomposition of $\C^N$.}
With $N = qp$ and the Kronecker skip $A = A_q \otimes I_p$, we decompose $\C^N = \C^q \otimes \C^p$. Every matrix $T \in M_N(\C)$ is viewed as a $q \times q$ array of $p \times p$ blocks:
\[
T = \bigl(T^{(\alpha\beta)}\bigr)_{\alpha,\beta=1}^q, \qquad T^{(\alpha\beta)} \in M_p(\C).
\]
The conditional expectation onto the subalgebra $\mcB = M_q(\C)$ is
\begin{equation}\label{eq:cond-exp}
E_\mcB[T] := \Bigl(\tfrac{1}{p}\Tr_p\, T^{(\alpha\beta)}\Bigr)_{\alpha,\beta=1}^q \;\in M_q(\C),
\end{equation}
which averages within each $p \times p$ sub-block while preserving the $q \times q$ block indices. The scalar trace factors as $\frac{1}{N}\Tr\, T = \frac{1}{q}\Tr_q\, E_\mcB[T]$.

\medskip
\noindent\textbf{Step 2: $\mcB$-valued order parameter.}
Define the $q \times q$ matrix-valued order parameter
\begin{equation}\label{eq:M-def}
M(z) := E_\mcB\!\bigl[Y(zI_N - Y^\top\! Y)^{-1}\bigr] = E_\mcB[G_{12}(z)] \;\in M_q(\C),
\end{equation}
where $G_{12}$ is the off-diagonal block of the $2N \times 2N$ Green's function~\eqref{eq:block-lin}. When $q = 1$, $E_\mcB$ reduces to $\frac{1}{N}\Tr$ and $M(z) = m(z)$.

\medskip
\noindent\textbf{Step 3: Wick contraction with Kronecker structure.}
We repeat the Schwinger--Dyson calculation of Section~\ref{app:scalar-recovery}, now applying $E_\mcB$ instead of $\frac{1}{N}\Tr$ to each block of~\eqref{eq:sd-identity}.

The random matrix $X = DW$ has i.i.d.\ $\mathcal{N}(0,\sigma^2/N)$ entries. The Wick contraction (Stein's lemma applied entry-by-entry as before) produces terms quadratic in the Green's function. The key structural point is that the contraction respects the $q \times q$ block decomposition:
\begin{itemize}
\item Within a given $p \times p$ sub-block $(\alpha,\beta)$, the sum over the $p^2$ entries of $X^{(\alpha\gamma)}$ paired with the resolvent produces the sub-block average $\frac{1}{p}\Tr_p(\cdot)$---exactly the $(\alpha,\gamma)$ entry of $E_\mcB[\cdot]$.
\item Across different sub-blocks, the independence of distinct entries of $X$ ensures no spurious cross-terms.
\end{itemize}
Concretely, the $(\alpha,\beta)$ component of the self-energy is:
\begin{equation}\label{eq:ov-wick}
\Sigma^{(\alpha\beta)} = \sigma^2 \sum_{\gamma=1}^q M_{\alpha\gamma}(z)\,\delta_{\gamma\beta}\, I_p
\;=\; \sigma^2\, M_{\alpha\beta}(z)\, I_p,
\end{equation}
where $M_{\alpha\beta}(z) = \frac{1}{p}\Tr_p\,(G_{12})^{(\alpha\beta)}$ is the $(\alpha,\beta)$ entry of the order parameter~\eqref{eq:M-def}. As a full $N \times N$ matrix, the self-energy is $\Sigma = \sigma^2(M(z) \otimes I_p)$: a Kronecker product that preserves the block structure.

\medskip
\noindent\textbf{Step 4: Dressed skip and $q \times q$ closure.}
The self-energy shifts the skip connection within the $q \times q$ block algebra:
\begin{equation}\label{eq:dressed-skip-ov}
A = A_q \otimes I_p \;\longrightarrow\; A_h(M) \otimes I_p, \qquad A_h(M) := A_q + \sigma^2 M(z) \;\in M_q(\C).
\end{equation}
Because $A_h(M) \otimes I_p$ is again Kronecker-structured, the dressed Green's function inherits the same block decomposition. Applying $E_\mcB$ to the Schur complement formula~\eqref{eq:G12} with $Y$ replaced by $A_h \otimes I_p$ collapses the $N \times N$ resolvent to a $q \times q$ one:
\begin{equation}\label{eq:kronecker-collapse-final}
M(z) = A_h(M)\,\bigl(zI_q - A_h(M)^\top A_h(M)\bigr)^{-1},
\end{equation}
which is the matrix Dyson equation~\eqref{eq:matrix-dyson}. The collapse works because $E_\mcB[(B \otimes I_p)(C \otimes I_p)^{-1}] = B\,C^{-1}$ for any $B, C \in M_q(\C)$ with $C$ invertible: the $p \times p$ identity factors cancel in the conditional expectation.

\medskip
\noindent\textbf{Step 5: Spectral density recovery.}
The scalar Cauchy transform of $Y^\top Y$ is
\[
G(z) = \frac{1}{q}\Tr_q\bigl(zI_q - A_h(M)^\top A_h(M)\bigr)^{-1},
\]
and the spectral density is again obtained by Stieltjes inversion~\eqref{eq:stieltjes-inversion}. The $q \times q$ equation~\eqref{eq:kronecker-collapse-final} is solved by Newton's method in $\C^{q^2}$ at cost $O(q^6)$ per step (Appendix~\ref{app:matrix-dyson-details}), compared with $O(N^3)$ for the full resolvent. Since $q \leq 4$ in practice, this is a dramatic reduction.

\begin{remark}[Verification: $q = 1$ recovery]\label{rem:verification_recovery}
When $q = 1$, $A_q = a \in \R$, $M(z) = m(z) \in \C$, and~\eqref{eq:kronecker-collapse-final} becomes
\[
m = \frac{a + \sigma^2 m}{z - (a + \sigma^2 m)^2},
\]
which is exactly~\eqref{eq:dyson-fp} evaluated for $A_h = (a + \sigma^2 m)\,I_N$.
\end{remark}

\section{Numerical Implementation Details}\label{app:numerical-details}

This appendix collects the implementation details for the numerical pipelines described in Section~\ref{sec:numerical}.

\subsection{Proof of the $z_1$-mapping (Proposition~\ref{prop:z1-mapping})}\label{app:z1-proof}

\begin{proof}
Denote the single-layer S-transform by $S_1(w)$.
By the definition of the S-transform (Definition~\ref{def:s-transform}), we have the identity
\begin{equation}\label{eq:s-cauchy-relation}
\frac{G_1(z)}{zG_1(z) - 1} = S_1\bigl(zG_1(z) - 1\bigr).
\end{equation}
Setting $w_1 := z_1 G_1(z_1) - 1$, this reads $S_1(w_1) = G_1(z_1) / w_1$.

For $L$ identical layers, the product Cauchy transform $G_L$ satisfies the same structural identity with $S_L = S_1^L$:
\begin{equation}
\frac{G_L(z_L)}{z_L G_L(z_L) - 1} = S_L\bigl(z_L G_L(z_L) - 1\bigr) = S_1\bigl(z_L G_L(z_L) - 1\bigr)^L.
\end{equation}

The key observation is that if we set $w_L := z_L G_L(z_L) - 1 = w_1$ (i.e., we demand that $z_1$ and $z_L$ correspond to the \emph{same} point $w$ in the S-transform domain), then
\begin{equation}
\frac{G_L(z_L)}{w_1} = S_1(w_1)^L = \left(\frac{G_1(z_1)}{w_1}\right)^L,
\end{equation}
so that $G_L(z_L) = w_1 \cdot (G_1(z_1)/w_1)^L = G_1(z_1)^L / w_1^{L-1}$.

From the constraint $z_L G_L(z_L) - 1 = w_1 = z_1 G_1(z_1) - 1$, we obtain
\begin{equation}
z_L = \frac{w_1 + 1}{G_L(z_L)} = \frac{z_1 G_1(z_1)}{G_L(z_L)}.
\end{equation}

Substituting $G_L(z_L) = G_1(z_1)^L / w_1^{L-1}$:
\begin{equation}
z_L = z_1 G_1(z_1) \cdot \frac{w_1^{L-1}}{G_1(z_1)^L} = z_1 \cdot \frac{w_1^{L-1}}{G_1(z_1)^{L-1}} = z_1 \left(\frac{w_1}{G_1(z_1)}\right)^{L-1},
\end{equation}
and since $w_1 = z_1 G_1(z_1) - 1$:
\begin{equation}
z_L = z_1 \left(\frac{z_1 G_1(z_1) - 1}{G_1(z_1)}\right)^{L-1},
\end{equation}
which is the first part of~\eqref{eq:z1-to-zL}. The relation $G_L(z_L) = z_1 G_1(z_1)/z_L$ then follows from $z_L G_L(z_L) = z_1 G_1(z_1)$.
\end{proof}

\subsection{Scalar Dyson solver: analytical derivative}\label{app:scalar-jacobian}

For each spectral parameter $z = x + i\eta$ with $\eta > 0$, we solve~\eqref{eq:dyson-fp} by Newton's method applied to $K(m, z) := \frac{1}{N}\Tr\bigl(A(m)(zI_N - A(m)^\top A(m))^{-1}\bigr) - m = 0$. Writing $R = zI_N - A_m^\top A_m$ where $A_m = A + \sigma^2 m I_N$, the analytical derivative is
\begin{equation}\label{eq:scalar-analytical-jacobian-app}
\frac{dK}{dm} = \frac{\sigma^2}{N}\Bigl[\Tr(R^{-1}) + \Tr\bigl(A_m R^{-1}(A_m^\top + A_m) R^{-1}\bigr)\Bigr] - 1,
\end{equation}
which reuses the already-computed $R^{-1}$ and avoids finite-difference overhead.

\subsection{Matrix-valued Dyson solver}\label{app:matrix-dyson-details}

When the skip connection has the Kronecker structure $A = A_q \otimes I_p$ (Section~\ref{sec:kronecker-skip}), the scalar Dyson equation~\eqref{eq:dyson-fp} is replaced by the $q \times q$ matrix equation~\eqref{eq:matrix-dyson}. We solve this by Newton's method in $\C^{q^2}$.

\medskip
\noindent\textbf{Vectorization.}
Writing $K(M, z) = A_h(zI_q - A_h^\top A_h)^{-1} - M$ as in~\eqref{eq:matrix-dyson}, the Newton update at iterate $M_k$ requires solving
\begin{equation}\label{eq:newton-matrix-dyson-app}
J_k \, \mathrm{vec}(\Delta M) = -\mathrm{vec}(K(M_k, z))
\end{equation}
for $\Delta M \in M_q(\C)$, where $J_k = \partial \mathrm{vec}(K) / \partial \mathrm{vec}(M) \in \C^{q^2 \times q^2}$ is the Jacobian.

\medskip
\noindent\textbf{Analytical Jacobian.}
Let $R = zI_q - A_h^\top A_h$. The Jacobian has the closed form
\begin{equation}\label{eq:analytical-jacobian-app}
J = \sigma^2\bigl(R^{-\top} \otimes I_q\bigr) + \sigma^2\bigl(R^{-\top}A_h^\top \otimes A_h R^{-1}\bigr)\Pi_q + \sigma^2\bigl(R^{-\top} \otimes A_h R^{-1} A_h^\top\bigr) - I_{q^2},
\end{equation}
where $\Pi_q \in \R^{q^2 \times q^2}$ is the \emph{commutation matrix} satisfying $\Pi_q \, \mathrm{vec}(X) = \mathrm{vec}(X^\top)$ for any $X \in M_q(\C)$, and $\otimes$ denotes the Kronecker product of matrices.

\medskip
\noindent\textbf{Damping and initial guesses.}
The Newton step is damped: if $\norm{\Delta M}_F > 10 \max(\norm{M_k}_F, 1)$, the step is rescaled to this bound. The iteration is seeded from multiple initial guesses---$M_0 = 0$, $M_0 = z^{-1}I_q$, $M_0 = (0.05 + 0.05i)I_q$---and the best-converging solution is selected. As in the scalar case, continuation from the previous $z$-point provides the primary seed.\label{rem:branch}

\medskip
\noindent\textbf{Matrix spectral parameter.}
The same Newton procedure extends to the generalized Dyson equation at a matrix spectral parameter $b \in M_q(\C)$:
\begin{equation}\label{eq:dyson-matrix-arg}
K(M, b) := A_h(M) \cdot \bigl(b - A_h(M)^\top A_h(M)\bigr)^{-1} - M = 0,
\end{equation}
by replacing $zI_q$ with $b$ throughout. This is needed for the Psi-inversion in the operator-valued S-transform (Section~\ref{sec:ov-pipeline-numerics}). The Jacobian~\eqref{eq:analytical-jacobian-app} generalizes by substituting $R = b - A_h^\top A_h$.

\subsection{Heterogeneous multi-layer: $w$-domain Newton}\label{app:w-domain}

For $L$ layers with different scalar skip matrices $A^l$ (the case $q = 1$), the S-transform product generalizes to $S_{J^\top J}(w) = \prod_{l=1}^L S_l(w)$, but the $z_1$-mapping (Proposition~\ref{prop:z1-mapping}) no longer applies. Instead, we solve for the common $w$-domain variable.

\medskip
\noindent\textbf{Outer Newton.}
For each target $z_L = x + i\eta$ on the evaluation grid, we seek $w \in \C$ such that
\begin{equation}\label{eq:w-domain-outer-app}
z_L(w) := \frac{w + 1}{w \cdot \prod_{l=1}^L S_l(w)} = z_L^{\mathrm{target}}.
\end{equation}
This is solved by Newton's method on $F(w) := z_L(w) - z_L^{\mathrm{target}} = 0$, with numerical derivative $dF/dw$.

\medskip
\noindent\textbf{Inner Newton.}
Each evaluation of $S_l(w)$ requires solving the scalar Psi-inversion $z_l G_l(z_l) - 1 = w$ for $z_l$, then $S_l(w) = G_l(z_l)/w$. The equation $z G(z) = w + 1$ has a spurious root at $z \to \infty$ where $zG(z) \to 1$; the Newton solve is therefore seeded at $z_l^{(0)} = z_L^{\mathrm{target}}$ (the physical branch) rather than the naive guess $(w+1) \cdot 1.05$.

\medskip
\noindent\textbf{Recovery and continuation.}
Once $w$ is found, $G_L(z_L) = (w + 1)/z_L$. All state---the outer $w$ and the inner $(z_l, m_l)$ for each layer---is carried forward to the next grid point, providing warm starts that maintain branch tracking across the $z$-grid. The outer step is damped: $\abs{\Delta w}$ is capped at $10 \max(\abs{w}, 1)$.

\subsection{Operator-valued pipeline: implementation}\label{app:ov-pipeline}

For $L > 1$ heterogeneous layers with $q > 1$, the full operator-valued pipeline involves a \emph{triple-nested Newton} structure.

\medskip
\noindent\textbf{Level 1: Outer Newton on $W \in M_q(\C)$.}
For each target $z = x + i\eta$, we solve the consistency equation~\eqref{eq:outer-consistency} by Newton's method on $q^2$ complex unknowns. The residual is
\begin{equation}\label{eq:ov-outer-residual-app}
F(W) = (W + I_q)\bigl(S_{\mathrm{prod}}(W) \cdot W\bigr)^{-1} - z I_q \in M_q(\C),
\end{equation}
with the Jacobian $\partial \mathrm{vec}(F)/\partial \mathrm{vec}(W) \in \C^{q^2 \times q^2}$ computed by $q^2$ finite-difference perturbations, each requiring a full evaluation of $S_{\mathrm{prod}}(W)$.

\medskip
\noindent\textbf{Level 2: Twisted S-product.}
Each evaluation of $S_{\mathrm{prod}}(W)$ computes the $L$-fold twisted recursion~\eqref{eq:twisted-fold}, calling $S_l^{(\mcB)}(W_{\mathrm{twisted}})$ once per layer.

\medskip
\noindent\textbf{Level 3: Psi-inversion per layer.}
Each $S_l^{(\mcB)}$ evaluation requires solving the Psi-inversion~\eqref{eq:psi-inversion} for $b_l \in M_q(\C)$, which itself invokes the matrix-argument Dyson solver (Section~\ref{app:matrix-dyson-details}).

\medskip
\noindent\textbf{Backtracking line search.}
The outer Newton step $\Delta W$ is accepted with step size $\alpha \in \{1, 0.5, 0.25, 0.1, 0.05, 0.02, 0.01\}$, choosing the largest $\alpha$ such that $\norm{F(W + \alpha \Delta W)}_F < \norm{F(W)}_F$ and all inner solves (Levels 2--3) converge.

\medskip
\noindent\textbf{Warm-starting and continuation.}
The per-layer Psi-inversion solutions $(b_l, M_l)$ are cached across both Newton iterations (within a single $z$-point) and across consecutive $z$-points, enabling effective continuation. The initial guess for $b_l$ is selected from a hierarchy:
\begin{enumerate}
\item \emph{Continuation hint}: the cached $b_l$ from the previous solve;
\item \emph{Scalar Newton}: find $z$ such that $z G_{\mathrm{scalar}}(z) = \frac{1}{q}\Tr(W + I_q)$, then set $b = zI_q$;
\item \emph{Per-sector diagonal hint} (when $A_q$ is diagonal): solve $q$ independent scalar Psi-inversion problems, one per diagonal entry of $A_q$;
\item \emph{Asymptotic fallback}: $b = W + I_q$.
\end{enumerate}

\subsection{Parallelism}\label{app:parallelism}

The evaluation grid $\{z_k = x_k + i\eta\}_{k=1}^K$ is split into $n_{\mathrm{jobs}}$ contiguous chunks for parallel processing.

\medskip
\noindent\textbf{Sequential pre-sweep.}
To initialize each chunk with a valid warm start, the first $z$-point of each chunk is solved \emph{sequentially} in order, carrying the full continuation state $(W, \{b_l, M_l\}_{l=1}^L)$ from chunk to chunk. Without this pre-sweep, cold-starting a chunk at a $z$-point deep inside the spectral support would require the full initial-guess hierarchy and may converge to the wrong branch.

\medskip
\noindent\textbf{Parallel dispatch.}
After the pre-sweep, each chunk is dispatched to a separate process via \texttt{ProcessPoolExecutor}. Thread-level parallelism (\texttt{ThreadPoolExecutor}) is avoided because the BLAS routines for small $q \times q$ matrices do not release the GIL.

\subsection{Monte Carlo validation}\label{app:mc-validation}

Theoretical predictions are validated by direct Monte Carlo sampling.

\medskip
\noindent\textbf{Single layer.}
For each sample $s = 1, \ldots, n_{\mathrm{MC}}$: draw $W^{(s)}$ with i.i.d.\ $\mathcal{N}(0, 1/N)$ entries, form
\begin{equation}
Y^{(s)} = (A_q \otimes I_p) + D^{(s)} W^{(s)},
\end{equation}
and compute the eigenvalues of $(Y^{(s)})^\top Y^{(s)}$ via a symmetric eigensolver. Here $D^{(s)} = \mathrm{diag}(\phi'(\sqrt{q^l} \, z_i))$ with $z_i \sim \mathcal{N}(0,1)$ for nonlinear activations, or $D^{(s)} = \sqrt{\sigma^2}\, I_N$ for the linear case.

\medskip
\noindent\textbf{Multi-layer.}
The Jacobian $J^{(s)} = \prod_{l=1}^L Y_l^{(s)}$ is formed by multiplying $L$ independent layer matrices, and the eigenvalues of $(J^{(s)})^\top J^{(s)}$ are computed. The empirical spectral density is obtained by histogramming the pooled eigenvalues across all samples with \texttt{density=True} normalization.

\subsection{Activation moments and signal propagation}\label{app:signal-prop}

The self-energy $\sigma^2$ entering the Dyson equation depends on the activation function $\phi$ and the signal propagation dynamics across layers.

\medskip
\noindent\textbf{Moment computation.}
The key activation moments are computed by Gauss--Hermite quadrature~\cite{golub1969calculation} (probabilists' Hermite polynomials, $n_{\mathrm{nodes}} = 100$):
\begin{equation}\label{eq:activation-moments-app}
\psi(v) = \E\bigl[\phi'(\sqrt{v}\, Z)^2\bigr], \quad \kappa(v) = \E\bigl[\phi(\sqrt{v}\, Z)^2\bigr], \quad Z \sim \mathcal{N}(0, 1).
\end{equation}

\medskip
\noindent\textbf{Signal propagation recursion.}
Given weight variance $\sigma_w^2$ and per-layer twist matrices $A_q^l$, the pre-activation variance $q^l$, self-energy $\sigma^2_l$, and signal variance $v^l$ propagate as
\begin{equation}\label{eq:signal-prop-app}
q^l = \sigma_w^2\,(v^{l-1} + \sigma^2_{l-1}), \quad \sigma^2_l = \psi(q^l), \quad v^l = \kappa(q^l) + \frac{\norm{A_q^l}_F^2}{q}\, v^{l-1}.
\end{equation}
The Dyson input is $\sigma^2_l = \sigma_w^2 \cdot \psi(q^l)$, and the critical dynamical isometry condition (Section~\ref{sec:dyson-resnet}) generalises to $L \sigma^2 = 1/2$ for identical layers.

%!TEX root = ../../main.tex

\section{Spectral Stalling: From Jacobian Spectrum to Optimiser Dynamics}\label{app:spectral-stalling}

The main text analyses the \emph{shape} of the Jacobian spectrum as a function of the skip-connection matrix $A_q$. This appendix addresses the complementary question: how does that spectral shape affect gradient-based training? We argue that four independent mechanisms each impose a \emph{hard cutoff} $\sigma^*$ on the singular-value spectrum, below which gradient directions are operationally invisible to the optimiser. The binding cutoff is
\begin{equation}\label{eq:sigma-star}
\sigma^* = \max\!\bigl(\sigma^*_{\mathrm{time}},\; \sigma^*_{\mathrm{Adam}},\; \sigma^*_{\mathrm{bf16}},\; \sigma^*_{\mathrm{noise}}\bigr).
\end{equation}
Below $\sigma^*$, the network trains in a subspace of effective dimension $n_{\mathrm{eff}} = n\!\int_{\sigma^*}^{\infty}\!d\mu_\Sigma(\sigma)$. We call this \emph{spectral stalling}, by analogy with aerodynamic stall: beyond a critical threshold, the gradient in a given direction no longer produces useful lift.

\subsection{Setup: spectral filtering of the backward pass}\label{app:stalling-setup}

The backward error signal at layer $l$ is $\delta_l = J_{l\to L}^\top \,\nabla_{\mathbf{h}_L}\!\mathcal{L}$, where $J_{l\to L} = \prod_{k=l+1}^{L} J_k$ with $J_k = D_k W_k + A_k$. Decomposing in the SVD $J_{l\to L} = U\Sigma V^\top$:
\begin{equation}\label{eq:spectral-filtering}
\delta_l = V\,\Sigma\,U^\top\,\nabla_{\mathbf{h}_L}\!\mathcal{L}, \qquad (\delta_l)_i = \sigma_i\,\alpha_i, \quad \alpha = U^\top\nabla\mathcal{L}.
\end{equation}
Each component of the loss gradient is multiplied by the corresponding singular value~$\sigma_i$. Directions with $\sigma_i \approx 0$ receive no gradient information. This is not the classical vanishing-gradient problem (which concerns the \emph{mean} of $\sigma$); it is a \emph{directional filtering} that depends on the full spectral measure~$\mu_\Sigma$.

\subsection{Cutoff I: exponential timescale separation}\label{app:cutoff-time}

Saxe, McClelland, and Ganguli~\cite{saxe2014exact} solved the gradient-flow dynamics of deep linear networks in closed form. In the SVD basis of the input--output correlation, each singular mode $i$ evolves with a sigmoidal trajectory whose time constant satisfies $\tau_i \propto 1/\sigma_i(0)$. Modes with small $\sigma_i$ take exponentially longer to converge; after a training budget of $T$ steps, only modes with $\tau_i \lesssim T$ have appreciably moved.

This yields a finite-time cutoff:
\begin{equation}\label{eq:sigma-time}
\sigma^*_{\mathrm{time}} \;\sim\; \frac{C}{T\,\eta\,\sigma_{\max}},
\end{equation}
where $C$ depends on the loss landscape curvature. Relative to the dominant mode, the number of steps to halve the error in direction $i$ scales as $t_{1/2}^{(i)} / t_{1/2}^{(\max)} = (\sigma_{\max}/\sigma_i)^2$~\cite{jacot2018ntk}. A mode with $\sigma_i = 0.01\,\sigma_{\max}$ requires $10^4\times$ more steps---it is simply not reached within any practical training budget.

Arora et al.~\cite{arora2019implicit} showed that gradient descent with small initialisation has an implicit bias toward low-rank solutions: $\dot{\sigma}_i \propto \sigma_i^{2-2/N}$, so large singular values accelerate while small ones stagnate. Pesme and Flammarion~\cite{pesme2023saddle} made this precise for diagonal linear networks: gradient flow sequentially jumps from saddle to saddle, each corresponding to the activation of a new singular mode, with jump times computable via a recursive algorithm.

\subsection{Cutoff II: Adam's $\epsilon$-threshold}\label{app:cutoff-adam}

Adam~\cite{kingma2015adam} maintains running estimates $m_i \approx \mathbb{E}[g_i]$ and $v_i \approx \mathbb{E}[g_i^2]$ of the gradient in each coordinate, with normalised update $m_i/(\sqrt{v_i} + \epsilon)$. Balles and Hennig~\cite{balles2018dissecting} showed this ratio is essentially the signal-to-noise ratio of $g_i$. In the regime $\sigma_i |\alpha_i| \ll \epsilon$, the update becomes $\sigma_i \alpha_i / \epsilon$, which vanishes linearly in $\sigma_i$:
\begin{equation}\label{eq:sigma-adam}
\sigma^*_{\mathrm{Adam}} \;\sim\; \frac{\epsilon}{|\bar\alpha|},
\end{equation}
where $\bar\alpha$ is the typical loss-gradient component. With the standard $\epsilon = 10^{-8}$, this cutoff is rarely binding in practice---the other floors kick in first.

A structural limitation prevents Adam from compensating for spectral filtering: Adam operates \emph{per-parameter} (in the weight-matrix coordinate system), but the spectral filtering operates in the \emph{Jacobian singular-vector basis}, which is a rotation of the parameter basis that changes with the input. Zhang et al.~\cite{zhang2024adam} demonstrated this directly: ViT training degrades under random global rotations of the parameter space, even though the optimisation landscape is unchanged. Adam's advantages are \emph{representation-dependent}, not universal~\cite{depavia2025rotations}.

\subsection{Cutoff III: low-precision arithmetic (bf16 swamping)}\label{app:cutoff-bf16}

BFloat16 has a relative precision of $u_{\mathrm{bf16}} = 2^{-8} \approx 3.9 \times 10^{-3}$. When accumulating a gradient update $\Delta w$ into a weight $w$, if $|\Delta w| < u_{\mathrm{bf16}} \cdot |w|$, the update is \emph{swamped}: $\mathrm{fl}(w + \Delta w) = w$ exactly. The gradient is lost to rounding. Since the gradient update in direction $i$ scales as $\eta \sigma_i |\alpha_i|$:
\begin{equation}\label{eq:sigma-bf16}
\sigma^*_{\mathrm{bf16}} \;\sim\; \frac{u_{\mathrm{bf16}}\,|w|}{\eta\,|\bar\alpha|}.
\end{equation}
With $\eta = 10^{-4}$, $|w| \sim 1$, $|\bar\alpha| \sim 1$: $\sigma^*_{\mathrm{bf16}} \sim 40$. This is an aggressive cutoff---in bf16 without loss scaling, only the dominant singular directions receive non-zero updates.

Loss scaling~\cite{micikevicius2018mixed} extends the dynamic range of the gradient but does not increase relative precision: the ratio $\sigma_i/\sigma_j$ is preserved after scaling. It helps with gradient \emph{underflow} but not with accumulation swamping. Stochastic rounding~\cite{gupta2015limited} softens the hard cutoff to a probabilistic one (preserving small updates in expectation), but the effective learning rate in low-$\sigma$ directions is still drastically reduced because stochastic-rounding variance overwhelms the signal~\cite{zamirai2020bfloat16}.

\subsection{Cutoff IV: stochastic noise floor}\label{app:cutoff-noise}

For mini-batch SGD with batch size $B$, the stochastic gradient in singular direction $i$ has mean $\sigma_i \bar\alpha_i$ and variance $\sigma_i^2 \operatorname{Var}(\alpha_i) / B$. The signal-to-noise ratio $\mathrm{SNR}_i = |\bar\alpha_i|\sqrt{B}/\mathrm{std}(\alpha_i)$ is independent of $\sigma_i$. However, the relative update magnitude in direction $i$ compared to the dominant direction is $\sigma_i / \sigma_{\max}$, regardless of noise. When $\sigma_i / \sigma_{\max} \ll 1$, the cumulative displacement in direction $i$ is negligible, giving:
\begin{equation}\label{eq:sigma-noise}
\sigma^*_{\mathrm{noise}} \;\sim\; \frac{\sigma_{\max}}{\mathrm{SNR}_{\max}} \cdot \frac{1}{\sqrt{B}}.
\end{equation}
This cutoff is alleviated by increasing batch size and is typically subdominant to the bf16 and timescale floors.

\subsection{The unified picture}\label{app:stalling-unified}

In typical modern training (Adam, bf16, batch size 256--4096, $T \sim 10^5$ steps), the four cutoffs have the following typical magnitudes:

\begin{center}
\begin{tabular}{lll}
\toprule
Cutoff & Expression & Typical order \\
\midrule
$\sigma^*_{\mathrm{time}}$ & $C / (T\eta\sigma_{\max})$ & $10^{-2}$--$10^{-1}$ \\
$\sigma^*_{\mathrm{Adam}}$ & $\epsilon / |\bar\alpha|$ & $10^{-4}$--$10^{-3}$ (rarely binding) \\
$\sigma^*_{\mathrm{bf16}}$ & $u_{\mathrm{bf16}}|w| / (\eta|\bar\alpha|)$ & $10^{-2}$--$10^{1}$ (\textbf{often dominant}) \\
$\sigma^*_{\mathrm{noise}}$ & $\sigma_{\max}/(\mathrm{SNR}_{\max}\sqrt{B})$ & $\sim 10^{-2}$ \\
\bottomrule
\end{tabular}
\end{center}

\noindent The bf16 swamping floor and the finite-time horizon are typically the binding constraints. This hierarchy means that spectral collapse in the Jacobian has direct, quantitative consequences for training: singular-value mass below $\sigma^*$ corresponds to capacity that the optimiser cannot access.

\subsection{Why only natural gradient could compensate}\label{app:natural-gradient}

The natural gradient update $\Delta\theta = -\eta\, F^{-1}\nabla_\theta\mathcal{L}$, where $F = \mathbb{E}[J_\theta^\top J_\theta]$ is the Fisher information (Gauss--Newton approximation), satisfies $F^{-1}\nabla\mathcal{L} = (\alpha_1, \ldots)$ in the Jacobian singular basis---the spectral filtering is \emph{exactly inverted}~\cite{amari1998natural}. Each direction receives a step proportional to $\alpha_i$ alone, independent of $\sigma_i$. This is why natural gradient converges at a rate independent of the condition number.

Adam approximates a \emph{diagonal} preconditioner in the \emph{parameter coordinate system}. The diagonal approximation discards the off-diagonal structure that couples parameters within a singular direction. Das et al.~\cite{das2024preconditioning} showed that Adam's effective condition number is $\mathcal{O}(\min(d, \kappa))$ for diagonal Hessians but offers no improvement for non-diagonal ones---precisely the case when spectral filtering operates in a rotated basis. K-FAC~\cite{martens2015kfac} and Shampoo~\cite{gupta2018shampoo} achieve partial spectral inversion via Kronecker-factored or per-dimension preconditioners, but neither fully inverts the Jacobian filtering. The hierarchy is: natural gradient (exact inversion) $>$ K-FAC/Shampoo (partial) $>$ Adam (coordinate-wise only) $>$ SGD (none).

\subsection{Connections to our spectral analysis}\label{app:stalling-connections}

\noindent\textbf{Computable stalling fraction.}
Our operator-valued Dyson pipeline gives the spectral measure $\mu_\Sigma$ as a function of the skip-connection matrix $A_q$, the weight variance $\sigma_w^2$, the activation function $\phi$, and depth $L$. Combined with the stalling threshold $\sigma^*$ from~\eqref{eq:sigma-star}, this yields a concrete, computable prediction for the stalling fraction:
\begin{equation}\label{eq:stalling-fraction}
\rho_{\mathrm{stall}}(A_q, L, \sigma^*) = \int_0^{\sigma^*} d\mu_\Sigma(\sigma).
\end{equation}

\noindent\textbf{Depth scaling.}
For doubly stochastic $A_q$ with interior eigenvalue $|\lambda_j| < 1$, the spectral measure develops a growing lower tail as $L$ increases, because the contribution of each contracting sector to the Jacobian decays as $|\lambda_j|^L$. The stalling fraction $\rho_{\mathrm{stall}}$ therefore increases with depth: \emph{deeper doubly stochastic networks stall in more directions}. For orthogonal $A_q \in O(q)$, all eigenvalues lie on the unit circle, the spectrum does not develop a lower tail, and $\rho_{\mathrm{stall}}$ remains small at all depths.

\medskip
\noindent\textbf{The bf16 interaction.}
The bf16 cutoff is particularly relevant for our setting because modern transformer training universally uses bf16 or fp16 mixed precision. Our spectral analysis can predict \emph{which} skip-connection structures produce spectra that fall below the bf16 swamping floor. If $\mu_\Sigma$ places significant mass in $[0, \sigma^*_{\mathrm{bf16}}]$, then bf16 training of that architecture will be effectively lower-dimensional---a prediction testable by comparing bf16 vs.\ fp32 training of networks with different skip-connection matrices.


\bibliographystyle{plain}
\bibliography{references}

\end{document}
